% SOURCE_AUTHORITY: sga5_fr_workpass.tex, lines 8126--8164 (LF-normalized unit).
% SOURCE_SHA256: 791F4EFFC5E02832D5D77ED03518C8156D6F07E4C8238B03545DB93D883FBB28
\begin{prooftext}
i) Para todo par \((m,r)\) de inteiros,
\[
\mathcal B_r^{(m)}(X)=\bigl(\mathcal B_r^{p,m-p}(X)\bigr)_{p\in\mathbb N}
\]
é um sub-\((\gr^\bullet(S),\mathcal D)\)-módulo graduado de \(\mathcal E_1^{(m)}(X)\). Em particular, para todo inteiro \(m\), os \(\mathcal B_r^{(m)}(X)\) formam uma sucessão crescente indexada por \(r\) de sub-\((\gr^\bullet(S),\mathcal D)\)-módulos graduados de \(\mathcal E_1^{(m)}(X)\). Usando que \(\mathcal E_1^{(n)}(X)\) e \(\mathcal E_1^{(n-1)}(X)\) são \((\gr^\bullet(S),\mathcal D)\)-módulos graduados noetherianos, deduz-se que as sucessões
\[
\bigl[\mathcal B_r^{(n)}(X)\bigr]_{r\in\mathbb N}
\qquad\text{e}\qquad
\bigl[\mathcal B_r^{(n-1)}(X)\bigr]_{r\in\mathbb N}
\]
São estacionárias.

Existe, portanto, um inteiro \(a\ge1\) tal que, para todo \((p,q)\), com \(p+q=n\) ou \(n-1\),
\begin{equation}
\mathcal B_r^{pq}(X)=\mathcal B_a^{pq}(X)\qquad(r\ge a).
\tag{A.3.1}
\end{equation}
Por conseguinte, para \(r\ge a+1\), as flechas \(d_r\) com origem e destino \(\mathcal E_r^{pq}(X)\) \((p+q=n)\) são nulas. Em outras palavras, os monomorfismos
\[
\mathcal E_s^{pq}(X)\hookrightarrow\mathcal E_r^{pq}(X)\qquad(s\ge r\ge a+1)
\]
deduzidos de (A.3.1) com \(p+q=n\) são, de fato, isomorfismos. Para \(s\ge r\ge a+1\), a comutatividade do diagrama \((d_3)\), com \(p+q=n\),
\[
\begin{array}{ccc}
\mathcal E_s^{pq}(X)&\xrightarrow[\epi]{\theta_s^{pq}}&E_\infty^{pq}(X_{p+s-1})\\
\wr\downarrow&&\downarrow\\
\mathcal E_r^{pq}(X)&\xrightarrow[\epi]{\theta_r^{pq}}&E_\infty^{pq}(X_{p+r-1})
\end{array}
\]
mostra que o morfismo de transição
\[
E_\infty^{pq}(X_{p+s-1})\longrightarrow E_\infty^{pq}(X_{p+r-1}),
\]
que, em qualquer caso, é um monomorfismo (cf. a)), é um epimorfismo e, portanto, um isomorfismo. Em outras palavras, o sistema projetivo
\[
r\longrightarrow\left[\bigoplus_p\gr^p T^n(X_{p+r})\right]_{r\ge1}
\]
é constante a partir de \(r=a\); por (4.1.3), \(T^n(X)\) verifica, portanto, (MLAR).
